%! Rational Functions and Holes — Unit 3, Lesson 3.4 — Warm-Up (Answer Key)
\documentclass[10pt]{article}
\usepackage{precalculus-article}
\usepackage{precalculus-key}
\newcommand{\TallMath}[1]{$\displaystyle #1\rule[-1.4em]{0pt}{3.2em}$}

\begin{document}

\pageheader{Unit 3, Lesson 3.4}{Warm-Up --- Answer Key}
\namedateperiod

\vspace{0.12in}

\begin{enumerate}[itemsep=14pt]

\item For the rational function $\displaystyle r(x) = \frac{5}{(x-3)(x+1)}$, find all
      vertical asymptotes. Write each as an equation.

      \vspace{0.1in}
      Vertical asymptote(s): \ans{$x = 3$ and $x = -1$}

      \textit{(Denominator zeros: $x-3=0 \Rightarrow x=3$ and $x+1=0 \Rightarrow x=-1$.
      Numerator is 5, nonzero at both. Both are VAs.)}

\item A table of values near $x = 2$ is shown below.

      \smallskip
      \renewcommand{\arraystretch}{1.3}
      \begin{tabular}{|c|c|c|c|c|c|}
      \hline
      $x$    & 1.9 & 1.99 & 2.001 & 2.01 & 2.1 \\
      \hline
      $r(x)$ & 47  & 497  & 501   & 51   & 5.7 \\
      \hline
      \end{tabular}
      \smallskip

      Based on the table, what does the behavior of $r(x)$ as $x \to 2$ indicate
      about $x = 2$? Circle the best answer and explain briefly.

      \medskip
      \textbf{A.} Hole at $x = 2$ \qquad \textbf{B.} \ans{Vertical asymptote at $x = 2$} \qquad \textbf{C.} Neither

      \vspace{0.1in}
      Explanation: \ans{The outputs grow without bound ($\to \infty$) as $x \to 2$ from both sides,
      indicating a vertical asymptote, not a hole. A hole would produce outputs approaching a
      finite value $L$.}

\item Factor the expression $x^2 + x - 6$ completely.

      \vspace{0.1in}
      $x^2 + x - 6 = $ \ans{$(x+3)(x-2)$}

\end{enumerate}

\end{document}
